\chapter{Towards real sensor data}
\label{ch:sensors}

\section{The sensor network}
A wireless sensor network (WSN) is being installed in the classroom. \todo{description of the WSN: number and
positions of the \cotwo{} sensors, sampling interval, accuracy, window and door contacts, wind sensors, people
counting; data format and access, as agreed with the lab}

\section{Data flow}
The model was trained on simulated \cotwo{} maps. To use it with the real sensors, the data has to flow from the
sensors to the model:
\begin{enumerate}
  \item the WSN measures the \cotwo{} at fixed positions, the window and door states and, if available, the number
        of people;
  \item the measurements are collected in a common format \todo{format agreed on};
  \item a preprocessing step brings them onto the model's time steps and converts the window states into the
        model inputs;
  \item the model forecasts the \cotwo{} for the next 3~minutes with an uncertainty;
  \item the forecast is shown or used to decide about the windows.
\end{enumerate}

\section{Open points}
\begin{itemize}
  \item \textbf{Few measurement points.} The model uses the \cotwo{} history of 900 points on a plane; the WSN
        gives only a few points. Since the network already works on point sets, it can be trained with only the
        sensor positions as history.
  \item \textbf{Window states instead of speeds.} The model needs the inflow speed of each window, while the WSN
        probably reports whether a window is open. A way to estimate the speeds is needed \todo{wind sensors?}.
  \item \textbf{All windows and doors.} The current dataset contains only W1, W3 and W6 with door~1 open and
        door~2 closed. In the real room all 8 windows and both doors can be opened. This needs more simulations and
        the door states as an additional model input.
\end{itemize}

\section{Synthetic sensor data}
Until the WSN delivers data, I generate synthetic sensor data from the simulations: the \cotwo{} at the sensor
positions every minute with realistic sensor noise, together with the window states and the number of people, in
the same format as the real data. With this data the whole chain from the sensors to the forecast can be built and
tested before the real data arrives. \todo{results}
